% Generated draft: requires mathematical and source review.
\documentclass[11pt]{article}
\usepackage[margin=25mm]{geometry}
\usepackage{fontspec}
\setmainfont{Latin Modern Roman}
\usepackage{amsmath,amssymb,mathtools,unicode-math}
\setmathfont{Latin Modern Math}
\usepackage{enumitem,needspace,xurl,hyperref}
\setlength{\parindent}{0pt}
\setlength{\parskip}{5pt}
\setlength{\emergencystretch}{3em}
\newcommand{\N}{\mathbb{N}}
\newcommand{\Z}{\mathbb{Z}}
\newcommand{\Q}{\mathbb{Q}}
\newcommand{\R}{\mathbb{R}}
\newcommand{\C}{\mathbb{C}}
\newcommand{\F}{\mathbb{F}}
\newcommand{\A}{\mathbb{A}}
\newcommand{\PP}{\mathbb{P}}
\newcommand{\E}{\mathbb{E}}
\newcommand{\eps}{\varepsilon}
\newcommand{\dd}{\,\mathrm{d}}
\begin{document}
\section*{UnsolvedMath ID: 11000173}
\textbf{Source problem code:} catalogue entry
\subsection{Definitions and mathematical statement}
Let $\mathcal M_g$ be the mapping class group of a closed orientable genus-$g$ surface and let
\[
 p:\mathcal M_g\longrightarrow\operatorname{Sp}(2g,\Z),\qquad
 \mathcal I_g=\ker p
\]
be the homological representation and Torelli group. A representation $\rho:\mathcal M_g\to A$ factors through $p$ when there is a homomorphism $\bar\rho:\operatorname{Sp}(2g,\Z)\to A$ with $\rho=\bar\rho\circ p$. Since $p$ is surjective in this setting, failure to factor is precisely
\[
 \mathcal I_g\not\subseteq\ker\rho,
 \quad\text{equivalently}\quad
 \exists t\in\mathcal I_g\text{ with }\rho(t)\ne1.
\]

\textbf{Problem 3.2.} Construct representations of $\mathcal M_g$, finite or infinite, which do not factor through $\operatorname{Sp}(2g,\Z)$. The preceding local source item 11000172 studies finite quotients and ends by motivating nonfinite quotients, supporting an image-size interpretation of ``finite or infinite.'' The present one-line request does not fix a target category, field, dimension, or a requirement of faithfulness. If linear representations are intended, one must specify $A=\operatorname{GL}(V)$ and whether dimension, rather than image size, is meant. Allowing completely unrestricted target groups would include the identity representation and would not express the intended constructive challenge, so the useful-target convention needs clarification.

The task is to detect some Torelli information, not necessarily to obtain an injective representation. The registry summary reports existing examples but also claims that mapping class groups are generally nonlinear and no faithful representations exist in large genus. Its own referenced local record 11000206 instead reports genus-two linearity and an unresolved genus-at-least-three linearity question. These status claims conflict; the nonlinearity assertion is not adopted as a theorem or as a limitation on this construction problem. The final inspirational paragraph introduces a following separate problem, not an additional condition here.
\subsection{Short English statement}
Construct useful representations that detect a nontrivial Torelli element and therefore cannot come solely from the integral symplectic action. Clarify whether finite/infinite refers to image or dimension; faithfulness is not required by the source.
\subsection{Sources}
\begin{itemize}
\item Local UnsolvedMath registry, permanent ID 11000173, source code catalogue entry; no proposer or proposal year supplied.
\item Farb, edited volume \emph{Problems on Mapping Class Groups and Related Topics}, Problem 3.2, PDF page 160; \url{https://www.math.uchicago.edu/~farb/papers/mcgbook.pdf}. The registry records PDF-font-structure extraction and access on 2026-07-29; the inventory distinguishes 329 labeled headings by source page.
\item Adjacent local ID 11000172 supplies the finite-quotient/nonfinite-quotient context. Local ID 11000206, explicitly referenced by the present research summary, reports a different linearity status, so the present summary's nonlinearity inference cannot be treated as verified.
\item No precise primary citation or explicit non-symplectic example is supplied in the present record.
\end{itemize}
\end{document}
